\subsection{کالبدشکافی ماژول \lr{smart\_weighted\_ensemble}}

\subsubsection*{۱. فلسفه طراحی، پارادایم معماری و تمثیل ملموس}
معماری \lr{smart\_weighted\_ensemble} رویکردی گروهی و ناهمگن است که خروجی چندین مکانیزم مختلف توجه (متراکم، خطی، کم‌رتبه و تنک) را از طریق یک شبکه گیتینگ هوشمند با ارزیابی عدم‌قطعیت تجمیع می‌نماید تا نقاط ضعف فردی هر مکانیزم توسط دیگری خنثی شود.

\textbf{تمثیل ملموس و مثال کاربردی:}
شورای عالی داوری متشکل از متخصصان و حقوق‌دانان زبده: دعاوی حقوقی پیچیده به جای ارجاع به یک قاضی منفرد، به هیئتی تخصصی عرضه می‌شود و رأی هر قاضی متناسب با درصد سوابق و احکام موفق پیشین او در آن موضوع ویژه، وزن‌دهی تحلیلی می‌گردد تا احتمال هرگونه خطای فردی به صفر میل کند.

\subsubsection*{۲. فرمول‌بندی و استنتاج عمیق ریاضی \lr{(Mathematical Derivation)}}
خروجی نهایی از ترکیب محدب خروجی شاخه‌های مختلف $\mathbf{Y}_k$ به دست می‌آید:
\begin{equation}
\mathbf{Y} = \sum_{k=1}^K w_k(\mathbf{X}) \mathbf{Y}_k, \quad \sum_{k=1}^K w_k(\mathbf{X}) = 1, \quad w_k(\mathbf{X}) \ge 0
\end{equation}
وزن‌های گیتینگ بر پایه اطمینان اطلاعاتی و انتروپی فعال‌سازی محاسبه می‌شوند:
\begin{equation}
\mathbf{w}(\mathbf{X}) = \operatorname{Softmax}\left( \mathbf{X} \mathbf{W}_{\text{gate}} + \mathbf{b}_{\text{gate}} \right)
\end{equation}
این ساختار، مقاومت و استحکام استثنایی در برابر نویز داده‌ها به ارمغان می‌آورد.

\subsubsection*{۳. کالبدشکافی الگوریتمی و جریان داده \lr{(Algorithmic Tensor Trace)}}
در پیاده‌سازی سورس پایتون در فایل ماژول \lr{smart\_weighted\_ensemble.py}، پردازش در متد \lr{forward} به صورت مرحله‌به‌مرحله و گام‌به‌گام صورت می‌پذیرد:
\begin{enumerate}
\item دریافت تانسور ورودی و ارسال به شبکه ارزیابی ضرایب گیتینگ هوشمند.
\item پردازش هم‌زمان ورودی در شاخه‌های ناهمگن توجه (خطی، تنک، اسپکترال).
\item محاسبه بردار وزن‌های تلفیق پویا متناسب با شرایط بافت متنی.
\item ضرب وزن‌دار خروجی شاخه‌ها با عملگر تجمیع تانسوری خطی.
\item ارسال خروجی تلفیقی به لایه پایانی بدون افزایش در بعد تانسور.
\end{enumerate}

\subsubsection*{۴. تحلیل همجوشی سخت‌افزاری و ترافیک حافظه \lr{(Hardware \& Roofline Analysis)}}
از دیدگاه مدل سقف کارایی (\lr{Roofline Model}) و معماری پردازنده‌های گرافیکی مدرن (مانند \lr{NVIDIA Ada / RTX 4070}):
\begin{itemize}
\item \textbf{همپوشانی محاسبات در جریان‌های موازی \lr{CUDA Streams}:} شاخه‌های مختلف به صورت هم‌زمان پردازش شده و هسته‌های بیکار را به کار می‌گیرند.
\item \textbf{کنترل بهینه حافظه:} شاخه‌های سبک بافر بسیار اندکی مصرف می‌کنند و اوج مصرف به 750.0 مگابایت محدود می‌شود.
\end{itemize}

\subsubsection*{۵. ارزیابی تجربی، بنچمارک و توجیه علمی نتایج}
بر اساس آزمون‌های ثبت‌شده در خط لوله ارزیابی پایان‌نامه:
\begin{itemize}
\item \textbf{وضعیت رسمی آزمون کیفی:} \lr{PASS}
\item \textbf{دقت ثبت‌شده روی بنچ‌مارک:} 90.14 درصد
\item \textbf{شاخص ارزیابی آماری (\lr{F1 Score}):} 87.7 درصد
\item \textbf{اوج مصرف حافظه گرافیکی (\lr{Peak VRAM}):} 750.0 مگابایت (در برابر ۴۲۰۰ مگابایت مدل مبنا)
\item \textbf{نرخ توان عملیاتی (\lr{Throughput}):} 32,500 توکن بر ثانیه
\end{itemize}

\textbf{توجیه علمی و علت‌شناسی جذاب نتایج:}
دستیابی به دقت کم‌نظیر 90.14 درصد و شاخص F1 معادل 87.7 درصد همراه با صرفه‌جویی 82.1 درصدی حافظه در برابر ۴۲۰۰ مگابایت مبنا، نشان‌دهنده استحکام بنیادین یادگیری گروهی در هم‌پوشانی خطاهای مکانیزم‌های خطی و متراکم است.

\subsubsection*{۶. دیاگرام معماری شکیل در لاتک \lr{(TikZ Architecture Flowchart)}}
جهت تبیین دقیق جریان داده تانسوری و نحوه تغییر شکل ماتریس‌ها، نمودار معماری این ماژول با استفاده از بسته قدرتمند \lr{TikZ} در محیط لاتک ترسیم شده است:

\begin{center}
\begin{tikzpicture}[scale=0.82, every node/.style={scale=0.82}, 
    node distance=1.1cm,
    auto,
    block/.style={rectangle, draw=blue!70, fill=blue!12, rounded corners=4pt, text centered, minimum height=0.85cm, font=\small},
    ens/.style={rectangle, draw=purple!70, fill=purple!15, rounded corners=4pt, text centered, minimum height=0.85cm, font=\small},
    outblock/.style={rectangle, draw=teal!80, fill=teal!15, rounded corners=4pt, text centered, minimum height=0.85cm, font=\small}
]
    \node [block] (input) {ورودی $\mathbf{X}$};
    \node [ens, right=1.0cm of input] (gate) {روتر عدم‌قطعیت};
    \node [ens, right=1.0cm of gate] (branches) {شاخه‌های ناهمگن};
    \node [outblock, right=1.0cm of branches] (out) {تلفیق هوشمند};

    \draw [->, thick, draw=blue!80] (input) -- (gate);
    \draw [->, thick, draw=purple!80] (gate) -- (branches);
    \draw [->, thick, draw=purple!80] (branches) -- (out);
\end{tikzpicture}
\end{center}

\subsubsection*{۷. پیوست تصویری و پرامپت تولید هوش مصنوعی}
برای ترسیم تجسم سه‌بعدی و گرافیک صنعتی این مکانیزم با ابزارهای هوش مصنوعی، پرامپت انگلیسی زیر با دقت تدوین شده است:

\vspace{0.2cm}
\begin{latin}
\footnotesize
\textbf{AI Technical Visualization Prompt:}\\
"A 3D visualization of the Smart Weighted Ensemble architecture. Multiple distinct crystalline streams of data uniting into a harmonious prism beam via a glowing intelligence nexus, vibrant golden and teal lighting, 8k resolution --ar 16:9"
\end{latin}

\begin{figure}[H]
    \centering
    \framebox[0.9\textwidth]{
    \begin{minipage}{0.85\textwidth}
    \vspace{1.8cm}
    \centering
    \textbf{[کادر الحاق تصویر گرافیکی سه‌بعدی ماژول \lr{smart\_weighted\_ensemble}]}\\[0.2cm]
    \small (بر اساس پرامپت فنی فوق تولید و در این کادر جایگزین می‌شود)
    \vspace{1.8cm}
    \end{minipage}
    }
    \caption{شماتیک و جریان داده تانسوری در ماژول \lr{smart\_weighted\_ensemble}}
    \label{fig:smart_weighted_ensemble}
\end{figure}

\vspace{0.5cm}
\hrule
\vspace{0.5cm}
